\documentclass[11pt]{article}
\usepackage[utf8]{inputenc}
\usepackage{amsmath,amssymb,amsfonts}
\usepackage[a4paper,margin=1in]{geometry}
\usepackage{booktabs}
\usepackage{hyperref}
\newcommand{\Mp}{M_{\rm Pl}}

\title{\textbf{A Dark-Matter-Coupled Dilaton Against Planck, ACT DR6, SPT-3G and DESI DR2:\\
an Apparent Phantom Crossing, the Hubble Constant,\\ and Six Pre-Registered Attempts to Raise It}}
\author{D.~L.~Eldridge\thanks{Independent researcher, Oklahoma, USA. dannyeldridge54@gmail.com.
Scripts, fits, locked rules and every number below: \url{https://github.com/dannyeldridge54/AEGIS}.}}
\date{11 November 2026}

\begin{document}
\sloppy
\maketitle

\begin{abstract}
We fit a dilaton dark-energy field coupled only to cold dark matter (the ``Spiral Twist'' flagship, FLAG-1: coupling
$\beta_{\rm eff}$ free, field-space slope fixed at $\kappa=1$) to Planck low-$\ell$, ACT DR6, SPT-3G D1, ACT DR6 lensing,
DESI DR2 BAO, a DES Y6 $S_8$ term and, in turn, Pantheon+, DES-Dovekie and Union3. Every test rule is fixed and hashed
before its fits are run. FLAG-1 improves on $\Lambda$CDM with the same $S_8$ term by $\Delta\chi^2=-4.71/-6.01/-7.69$
(one extra parameter), with $\beta_{\rm eff}\simeq0.04$ and $H_0\simeq67.2$ km\,s$^{-1}$Mpc$^{-1}$. Read with the
uncoupled-dark-matter assumption of the DESI analysis, the coupled field looks like dark energy whose $w$ crosses $-1$
near $z\simeq1$, although the field itself never crosses: a pre-registered DESI-style fit to a mock flagship universe
prefers $w_0w_a$ at $2.9\sigma$ with $w_0=-0.83$, $w_a=-0.45$, matching DESI DR2 + CMB + Pantheon+ ($2.8\sigma$). On the Hubble
constant, the supernova absolute magnitude implied by FLAG-1 ($M=-19.435$) differs from the SH0ES Cepheid calibration by
$0.186$ mag ($5.9\sigma$; $\Lambda$CDM $5.4\sigma$), but only by $0.083$ mag ($1.4\sigma$) from an approximate JWST
tip-of-the-red-giant-branch calibration. Six attempts to raise $H_0$ inside the model all fail, three of them as
pre-registered fits: an early $n=3$ field, a heavier electron before recombination, and extra dark radiation with
running. Three cheaper checks also fail: a universal coupling of the same dilaton to the electron breaks quasar and
atomic-clock limits by $\sim10^{5}$ and $\sim10^{4}$; a quadratic gravity coupling is not wanted by background data;
and the Pantheon+ Hubble-flow shows no local-void signature. Against the full DES Y3 cosmic-shear likelihood, FLAG-1
costs $\Delta\chi^2_{\rm WL}=+0.15/+1.33/+0.91$ relative to $\Lambda$CDM, within the pre-registered limit of $+2$ on all
three supernova sets; the KiDS-Legacy test is in progress. Failures are reported
at full size beside the passes.
\end{abstract}

\section{Introduction}
DESI DR2 BAO combined with the CMB and supernovae prefers evolving dark energy at $2.8$--$4.2\sigma$, with $w$
apparently crossing $-1$~\cite{desi2025}. A canonical scalar field cannot cross, but an analysis that assumes
uncoupled cold dark matter (CDM) can read a dark-energy--dark-matter energy exchange as a crossing. Dilaton fields that
couple to matter are natural in string-motivated models~\cite{damour1994,smith2025}. Here we test the simplest member of
this family, coupled to CDM only, against current data, and ask what it can and cannot do for the Hubble tension. The
work was run with pre-registration: each test's model, priors, data, gates and pass/fail bars were written down and
SHA-256 hashed before its fits ran. Deviations are logged before any rerun. AI assistance is stated in the
acknowledgements.

\section{Model and method}
The action is GR plus a scalar $\phi$ (reduced Planck units) with $V=\mu^4 e^{-\kappa\phi}$ ($\kappa=1$, fixed), and CDM
whose mass depends on $\phi$ with coupling $\beta_{\rm eff}$; baryons and photons are uncoupled. We use a modified
AxiCLASS (CLASS v3.3) with the coupling in the background and perturbations and a corrected HALOFIT treatment of the
field (the ``nl-scf-fix''). Fits use Cobaya minimisation, best of four starts. Data: Planck 2018 low-$\ell$ TT and EE,
ACT DR6 CMB-only with the Planck cut, SPT-3G D1 TT/TE/EE (lite), ACT DR6 lensing, DESI DR2 BAO, a DES Y6 $S_8$ Gaussian
term, and one supernova set per fit (Pantheon+, DES-Dovekie, Union3). The reference is $\Lambda$CDM fitted to the same
data with the same $S_8$ term (lower of two restarts). Rules: \texttt{field\_tests/rules/PREREG\_*.md} with their
\texttt{.sha256} files.

\section{The flagship fit}
\begin{center}
\begin{tabular}{lccc}
\toprule
 & Pantheon+ & DES-Dovekie & Union3\\
\midrule
$\Delta\chi^2$ vs $\Lambda$CDM+$S_8$ & $-4.71$ & $-6.01$ & $-7.69$\\
$\beta_{\rm eff}$ & 0.041 & 0.045 & 0.038\\
$H_0$ [km\,s$^{-1}$Mpc$^{-1}$] & 67.2 & 67.4 & 67.0\\
$S_8$ & 0.825 & 0.828 & 0.825\\
\bottomrule
\end{tabular}
\end{center}
The locked bar ($\Delta\chi^2\le-4$ on at least two of three sets) is met on all three. This test is not blind: the
model was chosen after earlier fits, and the rule says so. A closure detail of the build makes the reported $H_0$
0.108\% higher than the true $H(z{=}0)$ (about $0.07$ km\,s$^{-1}$Mpc$^{-1}$); $\chi^2$ values are unaffected.

\section{The DESI reading: an apparent phantom crossing}
Following the DESI convention, we define the dark energy that an uncoupled-CDM analysis would infer,
$\rho_{\rm DE,eff}(a)=\rho_{\rm tot}-\rho_b-\rho_r-\rho_\nu-\rho_{\rm cdm,0}a^{-3}$, with $w_{\rm eff}=-1-\frac13\,
d\ln\rho_{\rm DE,eff}/d\ln a$. Fitting $w_0+w_a(1-a)$ over $0<z<2$, weighted by $\rho_{\rm DE,eff}$, gives
$w_0=-0.847$ and $w_a=-0.33$ for all three fits. $w_{\rm eff}$ runs from $-1.09$ at $z=2$ through $-1.00$ at $z\simeq1$
to $-0.86$ today. The field's own equation of state stays above $-1$. Against DESI DR2 + CMB + Pantheon+ ($w_0=-0.838\pm
0.055$, $w_a=-0.62\pm0.21$), $w_0$ agrees and $w_a$ is $1.4\sigma$ weaker (one parameter at a time; the published
$w_0$--$w_a$ correlation is not used). The DES-Y5 and Union3 combinations sit $2$--$2.6\sigma$ away. This is a
consistency statement, not a prediction: FLAG-1 was itself fitted to DESI DR2 BAO. The mechanism is the same one
proposed for dark-axion and dark-baryon couplings~\cite{prl2026}.

\paragraph{A DESI-style fit of a flagship universe (pre-registered, DESI-MOCK).} We generate noiseless mock DESI DR2 BAO,
Pantheon+ and compressed CMB data ($100\theta_*$, $\omega_b$, $\omega_{bc}$ as seen at recombination) from the FLAG-1
Pantheon+ best fit, with the real covariances. We then fit $\Lambda$CDM and $w_0w_a$CDM with uncoupled CDM, as DESI does.
With BAO + CMB + SNe the mock prefers $w_0w_a$ by $\Delta\chi^2=11.1$ ($2.9\sigma$ for two parameters; DESI DR2 reports
$2.8\sigma$ for this combination). It recovers $w_0=-0.833$, $w_a=-0.45$, which is $0.1\sigma$ and $0.8\sigma$ from
DESI's values: the locked reading is REPRODUCES DESI'S SIGNAL. With BAO + CMB alone the mock gives only $1.0\sigma$
($w_0=-0.74$, $w_a=-0.70$), against DESI's $3.1\sigma$. In this model, then, much of the signal comes from the
supernovae. The CMB is compressed here, and FLAG-1 was fitted to the same DESI data, so this is a consistency test.

\section{The Hubble constant}
\paragraph{Inverse distance ladder.} At each best fit we compute the supernova distances and fit the Pantheon+ absolute
magnitude $M$ over the Hubble-flow SNe ($z>0.0233$, full STAT+SYS covariance). The SH0ES Cepheid hosts give
$M=-19.249\pm0.030$. FLAG-1 implies $M=-19.435\pm0.010$, a $0.186$ mag gap ($5.9\sigma$); $\Lambda$CDM+$S_8$ gives
$-19.421$ ($5.4\sigma$). On the Pantheon+ scale, the JWST/HST tip-of-the-red-giant-branch result ($H_0=70.39\pm1.94$,
\cite{freedman2025}) corresponds to $M\simeq-19.352\pm0.060$. The gap is then $0.083$ mag ($1.4\sigma$) for FLAG-1 and
$1.1\sigma$ for $\Lambda$CDM. This conversion is approximate: it rescales the SH0ES $M$ by the $H_0$ ratio.
\paragraph{No local void.} The same SH0ES-calibrated fit, with a separate $H_0$ in six redshift shells from $z=0.023$ to
$2.5$, is flat. The nearby shells ($z<0.15$) lie $0.5$--$0.7\%\pm1\%$ below the most distant one; a void that alone
explains the tension needs about $+8\%$. Freeing $H_0$ per shell improves $\chi^2$ by only $3.1$ for five extra
parameters, and a north/south split shows nothing. Whatever separates $73$ from $67$ is the same at every distance: it
sits in the calibration or before recombination, not in our neighbourhood.

\section{Six attempts to raise $H_0$ within the model}
\begin{center}
\small
\begin{tabular}{p{4.1cm}p{5.6cm}p{4.6cm}}
\toprule
Extension (rule) & Result & Reading\\
\midrule
Early $n=3$ field + dilaton (TWOFIELD) & $\Delta\chi^2$ vs FLAG-1 $-1.06/-1.00/+0.62$; $H_0$ $67.8/67.5/67.6$ & NOT WORTH IT (3/3); $H_0$ does not rise\\
Electron mass before $z=50$ (TWIST-ME) & $m_e/m_{e0}=1.009/0.998/1.002$; $\Delta\chi^2$ $-0.55/-0.33/-0.15$ & NOT WORTH IT (3/3)\\
Dark radiation + running (TWIST-DR) & $\Delta N_{\rm eff}\to0.000$; $\Delta\chi^2$ $-1.04/-0.47$ (Pantheon+ fit died) & NOT WORTH IT (2/3)\\
Universal coupling $m_e\propto e^{g\phi}$ & $\phi$ moves $0.12$ since recombination but $0.21$ since $z=0.89$ & breaks quasar and clock limits by $\sim10^5$, $\sim10^4$\\
Quadratic gravity coupling $G_{\rm eff}=1+A(\phi-\phi_0)^2$ (toy) & best $A=0.047$, $\Delta\chi^2=-0.08$; forcing an 8\% bump costs $+71$ & not wanted (background only)\\
Local void & shells flat (above) & no\\
\bottomrule
\end{tabular}
\end{center}
The first three are pre-registered fits on the full likelihood; the last three are exploratory checks, marked so in the
repository. In TWOFIELD a HALOFIT bug (the new field was treated as matter) made the locked engine gate fail. It was
logged as deviation D1, fixed, and the gate rerun once and passed (lensed $C_\ell$ then match the flagship to
$4\times10^{-5}$). In TWIST-ME, where $m_e$ rose the coupling fell ($\beta_{\rm eff}$ $0.041\to0.028$), so the two
partly do the same job. In the gravity toy, $(\phi-\phi_0)^2$ from the flagship's own field history peaks near $z\simeq2$,
the shape reconstructed in Ref.~\cite{scalartensor2025}. That analysis's $H_0$ gain relies on SH0ES-calibrated
supernovae and a free-form $G(z)$, and does not survive here. Our results match the literature's finding that, once DESI
BAO is included, recombination-era fixes cannot raise $H_0$ much without lowering $\Omega_m$ below what BAO
allows~\cite{recomb2026}.

\section{Weak lensing}\label{sec:wl}
A late-time variant with a dark-matter angle (S8-A) passed a one-number $S_8$ test but missed the full DES Y3 shear
($\Delta\chi^2_{\rm WL}=+1.76/+2.32/+2.76$) and KiDS-Legacy ($+2.96/+10.11/+11.30$) tests, so the flagship stayed the
base. A joint refit with DES Y3 shear inside the likelihood (rule JOINT-SHEAR) needed a quadratic emulator of the shear
$\chi^2$. It missed held-out points by up to $3.8$ (S8-A) and $2.2$ ($\Lambda$CDM); after the one allowed enlargement of
the calibration it still missed by up to $1.8$ and $2.0$, against a limit of $0.5$. No fits ran and no verdict exists.

The locked full-shear test (rule S8-WL, Part 2) was then run at the FLAG-1 best fits: cosmology fixed at each fit, the DES
Y3 nuisance parameters (intrinsic alignments, photometric-redshift shifts, shear calibration) minimised, best of seven
starts, against the matched $\Lambda$CDM fits treated the same way. FLAG-1 gives $\Delta\chi^2_{\rm WL}=+0.15$ (Pantheon+),
$+1.33$ (DES-Dovekie) and $+0.91$ (Union3): \textbf{PASSES} on all three (limit $+2$). The earlier Spiral fits without the
$S_8$ term were OVER ($+2.57/+3.47/+3.99$). The same test against KiDS-Legacy is running and will be reported whatever
it gives; weak lensing is settled for FLAG-1 only once both surveys are in.

\section{Relation to the axiodilaton}
Ref.~\cite{smith2025} couples a dilaton to all matter. The electron mass then varies through recombination, giving
$H_0\simeq69.2$ without SH0ES and an apparent phantom crossing, but the coupling needs screening to pass solar-system
tests. FLAG-1 is the other side of the same trade-off. Coupling to CDM only keeps it clear of laboratory, quasar and
solar-system limits and keeps the apparent crossing, but leaves $H_0$ near 67. Our universal-coupling check shows why
the Twist cannot simply take the other side: its field moves more after $z\simeq1$ than before recombination.

\section{Frozen predictions}\label{sec:pred}
Before the relevant data are public, the FLAG-1 best fits and the numbers below are frozen by SHA-256 under locked rules
in the archive. No refit is allowed under any of them. Each line is scored once, whatever the outcome.

\textbf{DESI DR3} (rule \texttt{PREREG\_PRED\_DESI\_DR3}; Pantheon+ fit primary). The fit gives $r_d=147.45$\,Mpc,
$H_0=67.16$, $\sigma_8=0.807$, $S_8=0.825$, and $f\sigma_8=0.419$, $0.457$, $0.459$, $0.450$, $0.425$, $0.373$ at
$z=0$, $0.3$, $0.5$, $0.7$, $1.0$, $1.5$.
\begin{itemize}\setlength\itemsep{0pt}
\item[P1] \emph{Dark-energy reading.} DESI's own $w_0w_a$ fit to DR3 BAO + CMB + Pantheon+ should land near
$w_0=-0.833$, $w_a=-0.451$. It is CONSISTENT if each value is within $2\sigma$, and FAILS if either is beyond $3\sigma$.
\item[P2] \emph{The signal persists.} The model FAILS if DR3 + CMB + Pantheon+ prefers $w_0w_a$ over $\Lambda$CDM at
less than $1\sigma$.
\item[P3] \emph{BAO distances.} $\chi^2$ of the DR3 points against the frozen $D_M/r_d$, $D_H/r_d$, $D_V/r_d$, with no
free parameter. It FAILS if $p<0.001$ and is CONSISTENT if $p\ge0.01$.
\item[P4] \emph{Growth.} The same test for DESI full-shape $f\sigma_8$ against the values above.
\end{itemize}

\textbf{Weak lensing} (rule \texttt{PREREG\_LENSING\_FORECASTS} with its FLAG-1 addendum). The forecasts are
$S_8=0.825/0.828/0.825$ (Pantheon+/DES-Dovekie/Union3), against $0.814/0.816/0.814$ for the matched $\Lambda$CDM fits. The
targets are Euclid DR1 and Rubin/LSST Y1 cosmic shear, plus the DES Y6 full-shear likelihood. The tests are the same
summary and full-shear tests as Section~\ref{sec:wl}, with the same $+2$ limit.

\textbf{Nulls} (rule \texttt{PREREG\_NULLS}). Because only dark matter couples, FLAG-1 predicts no effect in each of the
following:
\begin{itemize}\setlength\itemsep{0pt}
\item gravitational waves travel at $c$ with GR's friction ($\Xi_0=1$), and there is no circular polarisation of the GW
background;
\item no drift of the fine-structure constant;
\item no fifth force on ordinary matter;
\item $N_{\rm eff}=3.044$ at BBN;
\item cosmic birefringence $\simeq0^\circ$.
\end{itemize}
Any measurement of one of these effects at $\ge5\sigma$ fails that line.

\section{Limitations and how to falsify it}
FLAG-1 was not chosen blind. $\kappa=1$ is fixed by hand. The $S_8$ term is a Gaussian summary, not the full shear. The
JWST conversion is approximate. The $w_0$--$w_a$ comparison ignores the published correlation. The gravity check is
background only. The model is falsified by any FAIL among the frozen predictions above: chiefly if DESI DR3 + CMB drive $w_a$ to zero
while the coupling stays required, if growth or BAO data depart from the frozen values, or if the KiDS-Legacy shear
test at its best fits, or a future survey's, comes out OVER.

\section{Conclusions}
A dilaton coupled only to dark matter fits the newest CMB, BAO and supernova data better than $\Lambda$CDM. It reproduces
DESI's apparent phantom crossing without a phantom field, and passes laboratory tests by construction. It does not
resolve the gap to SH0ES; nothing we added inside it does. Its $H_0$ is within $1.4\sigma$ of the JWST red-giant
calibration. It passes the full DES Y3 shear test; the KiDS-Legacy test is pending.

\section*{Acknowledgements}
We thank the Planck, ACT, SPT, DESI, DES, Pantheon+ and Union teams for public data, and the developers of CLASS,
AxiCLASS, CAMB and Cobaya. The author used AI assistants (Claude, Anthropic) for code, running the fits and drafting;
every number is reproduced by archived scripts and locked rules, and the author takes full responsibility.

\begin{thebibliography}{99}
\bibitem{desi2025} DESI Collaboration, arXiv:2503.14738.
\bibitem{damour1994} T.~Damour, A.~M.~Polyakov, Nucl.\ Phys.\ B 423, 532 (1994), hep-th/9401069.
\bibitem{smith2025} A.~Smith, M.~Mylova, C.~van de Bruck, C.~P.~Burgess, E.~Di Valentino, arXiv:2512.13544.
\bibitem{prl2026} \emph{Apparent $w<-1$ and a lower $S_8$ from dark axion and dark baryons interactions}, Phys.\ Rev.\ Lett., doi:10.1103/w4qb-plk8.
\bibitem{freedman2025} W.~L.~Freedman et al., arXiv:2408.06153.
\bibitem{scalartensor2025} arXiv:2511.04610.
\bibitem{recomb2026} arXiv:2606.06495.
\bibitem{seto2025} O.~Seto et al., arXiv:2508.09025.
\bibitem{kanekar2015} N.~Kanekar et al., quasar methanol bound on $\Delta\mu/\mu$ (2015).
\bibitem{nakashima2008} M.~Nakashima et al., arXiv:0803.0247.
\end{thebibliography}
\end{document}
